空间散射矩阵的时间演化由演化 Lax 方程确定.
由于 $V$ 在 $k=a$ 处具有极点, 空间归一化的 Jost 解满足带右乘修正项的时间方程.
由此可分别求出连续与离散散射数据的演化, 并在共振原像处利用介质方程确定相容条件.

\begin{proposition}\label{prop:evolution}
设三场满足基本条件及规范 HMB 系统. 对避开共振原像的普通连续谱,
\begin{equation}\label{eq:chi-time}
\chi_{\pm,t}=V\chi_\pm-\chi_\pm\ii\lambda F\sg.
\end{equation}
散射矩阵满足 $S_x=0$ 和 $S_t=[\ii\lambda F\sg,S]$, 因而
\begin{equation}\label{eq:timeS}
S(z,t)=e^{\ii\lambda Ft\sg}S(z,0)e^{-\ii\lambda Ft\sg}.
\end{equation}
透射系数及其解析零点不随 $t$ 变化; 反射系数满足
\begin{equation}\label{eq:timer}
\rho(z,t)=e^{-2\ii\lambda Ft}\rho(z,0),\qquad
\widetilde\rho(z,t)=e^{2\ii\lambda Ft}\widetilde\rho(z,0).
\end{equation}
\end{proposition}
\begin{proof}
基本条件中的 $C^1(I;L^1)$ 正则性允许对相应 Jost 积分方程求时间导数.
在避开 $0,\pm\ii q_0,z_a^\pm$ 的谱参数紧集上, Volterra 核与时间导数的迭代由可积函数一致控制.
零曲率关系给出
\begin{equation*}
\partial_x(\chi_{\pm,t}-V\chi_\pm)=U(\chi_{\pm,t}-V\chi_\pm).
\end{equation*}
在各自归一化端, $V-V_\infty\to0$, 且
$V_\infty Y=\ii\lambda F Y\sg$.
因此 $\chi_{\pm,t}-V\chi_\pm$ 与 $-\chi_\pm\ii\lambda F\sg$ 满足同一空间方程和同一端态.
去除空间指数后应用 Volterra 唯一性, 得到 \eqref{eq:chi-time}.
同样的逐列论证适用于各解析列.

对散射关系求导, 有
\begin{equation*}
\chi_-S_t
=(V\chi_+-\chi_+\ii\lambda F\sg)
-(V\chi_--\chi_-\ii\lambda F\sg)S
=\chi_-[\ii\lambda F\sg,S].
\end{equation*}
连续谱上的 $\chi_-$ 可逆, 故得到 $S_t$ 方程及 \eqref{eq:timeS}.
将右端的对角矩阵相乘,
\begin{equation*}
S(z,t)=
\begin{pmatrix}
s_{11}(z,0)&e^{2\ii\lambda Ft}s_{12}(z,0)\\
e^{-2\ii\lambda Ft}s_{21}(z,0)&s_{22}(z,0)
\end{pmatrix}.
\end{equation*}
对角元不变, 非对角元除以相应透射系数即得 \eqref{eq:timer}.
透射系数的解析延拓由相应 Jost 列的同一论证确定, 所以离散零点不随 $t$ 变化.
\end{proof}

\begin{remark}
在 $k\ne a$ 处, 令
\begin{equation}\label{eq:full-jost}
\phi_\pm=\chi_\pm e^{\ii\lambda Ft\sg}
=\mu_\pm e^{\ii\theta\sg}.
\end{equation}
则 $\phi_{\pm,x}=U\phi_\pm$, $\phi_{\pm,t}=V\phi_\pm$.
\end{remark}

\begin{proposition}\label{prop:norming-evolution}
在命题 \ref{prop:evolution} 的条件下, 设 $\zeta_j$ 是简单离散零点.
按 \eqref{eq:norming-definition} 定义的空间范数数据满足
\begin{equation}\label{eq:norming-time}
b_j(t)=b_j(0)e^{2\ii\lambda(\zeta_j)F(k(\zeta_j))t},\qquad
c_j(t)=c_j(0)e^{2\ii\lambda(\zeta_j)F(k(\zeta_j))t}.
\end{equation}
相应留数系数为
\begin{equation}\label{eq:C-timephase}
C_j(x,t)=c_j(t)e^{-2\ii\lambda(\zeta_j)x}
=c_j(0)e^{2\ii\theta(x,t,\zeta_j)}.
\end{equation}
\end{proposition}
\begin{proof}
修正时间方程的两列分别满足
\begin{equation*}
\chi_{\pm,1,t}=V\chi_{\pm,1}-\ii\lambda F\chi_{\pm,1},\qquad
\chi_{\pm,2,t}=V\chi_{\pm,2}+\ii\lambda F\chi_{\pm,2}.
\end{equation*}
把 $\chi_{+,2}=b_j(t)\chi_{-,1}$ 求时间导数,
消去相同的 $V\chi_{-,1}$ 项, 得
\begin{equation*}
b_j'(t)\chi_{-,1}
=2\ii\lambda(\zeta_j)F(k(\zeta_j))b_j(t)\chi_{-,1}.
\end{equation*}
该列非零, 从而得到 $b_j$ 的标量常微分方程.
积分即得第一式; $s_{22}'(\zeta_j)$ 不随时间变化, 得到第二式.
再代入空间去相位因子, 得 \eqref{eq:C-timephase}.
\end{proof}

介质矩阵的共同端态为
\begin{equation}\label{eq:Rbg}
R_\infty=\begin{pmatrix}
\eta_0&\dfrac{\ii\eta_0q_0}{a}\\[6pt]
-\dfrac{\ii\eta_0q_0}{a}&-\eta_0
\end{pmatrix}
=\dfrac{\ii\eta_0}{a}U_\infty(a).
\end{equation}

\begin{proposition}\label{prop:resonance}
固定 $t$. 设三场满足基本条件及两条介质方程. 则在每个共振原像,
\begin{equation}\label{eq:resonance-conjugation}
\chi_\pm(z_a^\pm)^{-1}R\chi_\pm(z_a^\pm)
=\dfrac{\eta_0\lambda(z_a^\pm)}{a}\sg.
\end{equation}
散射矩阵满足
\begin{equation}\label{eq:resonance-reflection}
[S(z_a^\pm,t),\sg]=0,\qquad
\rho(z_a^\pm,t)=\widetilde\rho(z_a^\pm,t)=0.
\end{equation}
\end{proposition}
\begin{proof}
空间 Jost 构造在 $z_a^\pm$ 正则, 且 $\det\chi_\pm=\gamma(z_a^\pm)>0$.
将介质方程写为 $R_x=[U(a),R]$, 则
\begin{equation*}
\partial_x(\chi_\pm^{-1}R\chi_\pm)
=\chi_\pm^{-1}(R_x-[U(a),R])\chi_\pm=0.
\end{equation*}
由背景恒等式,
\begin{equation*}
Y(z_a)^{-1}R_\infty Y(z_a)=\dfrac{\eta_0\lambda(z_a)}{a}\sg.
\end{equation*}
$\lambda(z_a)=\pm\sqrt{a^2+q_0^2}$ 为非零实数.
在 $\chi_\pm$ 各自的归一化端, 振荡矩阵及其逆有界, 并与 $\sg$ 对易.
取极限就得到 \eqref{eq:resonance-conjugation}.
再使用 $\chi_+=\chi_-S$, 得
\begin{equation*}
S(z_a,t)^{-1}\dfrac{\eta_0\lambda(z_a)}{a}\sg S(z_a,t)
=\dfrac{\eta_0\lambda(z_a)}{a}\sg.
\end{equation*}
因系数非零, $S$ 必须对角; $\det S=1$ 保证两个对角元均非零, 故反射系数为零.
\end{proof}
\begin{remark}
式 \eqref{eq:resonance-conjugation} 给出
\begin{equation*}
R\chi_\pm(z_a)-\chi_\pm(z_a)\dfrac{\eta_0\lambda(z_a)}{a}\sg=0.
\end{equation*}
故修正时间方程 \eqref{eq:chi-time} 的共振留数相消.
该共轭恒等式同时给出介质矩阵的谱表示, 是三场重构的依据.
\end{remark}
